\section{Carry input history across blocks}
The earlier kernel accepts at most sixty-four output samples per call. Its input window
contains seven chronological history samples followed by that call's new samples.
For a file wrapper, form such a window for each channel, invoke the kernel, and retain
the last seven \emph{input} samples for the next block. The history does not contain
the newly filtered outputs. Using outputs would introduce feedback and change the
mathematical operation.

Suppose the first block has three frames. Its window begins with seven zeros and then
three new values. The next history is the last seven values of that ten-value window:
four zeros followed by the three actual input codes. Keeping only the three new codes
without their preceding zeros would leave four uninitialized history positions. The
same rule works for a one-frame block and for a full sixty-four-frame block.

At every split point, concatenated block results must equal a whole-file result with
the same beginning condition. Check block sizes one, three, seven and sixty-four.
A difference at a block boundary usually identifies wrong retained history, a changed
channel stride or a repeated rounding operation. Reset both channel histories at the
start of each new file, including after an earlier file operation failed.
